\documentclass[11pt]{article}
\usepackage[a4paper,margin=2cm]{geometry}
\usepackage{amsmath,amssymb}
\usepackage[hidelinks]{hyperref}
\allowdisplaybreaks
\title{Qpu.Processing.Verify.Receipt}
\author{Tsvetan Rouschev}
\date{@uuidna/qpu v1.1.1, CC-BY-NC-ND-4.0}
\begin{document}
\maketitle
\noindent\url{https://github.com/uuidna/qpu/blob/v1.1.1/lean/Qpu/Processing/Verify/Receipt.lean}

\paragraph{\texttt{Qpu.Verify.Receipt.address\_cheaper\_once}}
Even the first run is cheaper: lookup beats evaluation.
\[ \mathrm{Qpu.Verify.Receipt.sealed}\ 1\ \mathrm{Qpu.Verify.Receipt.checkDs}\ 0 < \mathrm{Qpu.Verify.Receipt.rerun}\ 1\ \mathrm{Qpu.Verify.Receipt.evalDs} \]
\noindent axioms: none; proof: \url{https://github.com/uuidna/qpu/blob/v1.1.1/lean/Qpu/Processing/Verify/Receipt.lean\#L23}

\paragraph{\texttt{Qpu.Verify.Receipt.breaks\_even\_by\_run\_two}}
Break-even arrives by run two: caching is already net-positive on the first repeat.
\[ \mathrm{Qpu.Verify.Receipt.sealed}\ 2\ \mathrm{Qpu.Verify.Receipt.checkDs}\ 0 < \mathrm{Qpu.Verify.Receipt.rerun}\ 2\ \mathrm{Qpu.Verify.Receipt.evalDs} \]
\noindent axioms: none; proof: \url{https://github.com/uuidna/qpu/blob/v1.1.1/lean/Qpu/Processing/Verify/Receipt.lean\#L27}

\paragraph{\texttt{Qpu.Verify.Receipt.caching\_has\_a\_cost}}
Control: caching is itself not free — the address check has nonzero cost.
\[ 0 < \mathrm{Qpu.Verify.Receipt.checkDs} \]
\noindent axioms: none; proof: \url{https://github.com/uuidna/qpu/blob/v1.1.1/lean/Qpu/Processing/Verify/Receipt.lean\#L34}

\paragraph{\texttt{Qpu.Verify.Receipt.fold\_is\_order\_sensitive}}
Control: the fold is order-dependent, so swapping two leaves changes the root — a tamper-evident receipt cannot commute silently.
\[ \mathrm{Qpu.Verify.Receipt.merkleFold} [1, 2] \ne \mathrm{Qpu.Verify.Receipt.merkleFold} [2, 1] \]
\noindent axioms: none; proof: \url{https://github.com/uuidna/qpu/blob/v1.1.1/lean/Qpu/Processing/Verify/Receipt.lean\#L44}

\paragraph{\texttt{Qpu.Verify.Receipt.speedup\_is\_thirtyfold}}
\~{}30× speedup: 15·30 ≤ 459 < 15·31, so the ledger factor is bracketed between 30 and 31 for the reference bundle.
\[ 30 \cdot \mathrm{Qpu.Verify.Receipt.checkDs} \le \mathrm{Qpu.Verify.Receipt.evalDs} \land \mathrm{Qpu.Verify.Receipt.evalDs} < 31 \cdot \mathrm{Qpu.Verify.Receipt.checkDs} \]
\noindent axioms: none; proof: \url{https://github.com/uuidna/qpu/blob/v1.1.1/lean/Qpu/Processing/Verify/Receipt.lean\#L31}
\end{document}
